\section{Topological spaces: open tests, coefficients, and support}

In this section $\mathrm{Top}$ denotes all small topological spaces; in
particular we do not impose the Hausdorff condition on a space over a finite
Alexandrov base. We write $\Pi X=\operatorname{Sing}(X)\in\mathcal S$ for its
weak homotopy type. An infinity-categorical localization is used whenever
the target is an infinity-category. Its ordinary homotopy category is the
ordinary localization occurring in the manuscript.

\begin{claim}{D.1}{Known}{The actual comparison supplied by open tests}
For a space $Y$, put
\[
 \begin{split}
 W^{\mathrm{opn}}_Y=\{u:X\longrightarrow X'\text{ over }Y:\quad&
   \Pi(u^{-1}U):\Pi(X\times_YU)\longrightarrow\Pi(X'\times_YU)\\
 &\text{is an equivalence for every open }U\subseteq Y\}.
 \end{split}
\]
Here $u^{-1}U$ means the displayed restricted map, not an inverse-image
operation on morphisms. This class contains the isomorphisms and satisfies
two-out-of-three and retract closure. There is a canonical functor
\[
 C_Y:(\mathrm{Top}/Y)[(W^{\mathrm{opn}}_Y)^{-1}]
       \longrightarrow\operatorname{CoShv}(Y;\mathcal S),\qquad
 C_Y(X)(U)=\Pi(X\times_YU).
\tag{D.1}
\]
Its values are hypercomplete cosheaves. This statement does not assert
that (D.1) is an equivalence for every $Y$.
\end{claim}
\begin{proof}
For each open $U$, inverse image gives a functor from the slice to
$\mathrm{Top}$. The asserted closure properties follow by taking the
intersection of the inverse images of the corresponding properties of
weak homotopy equivalences. The higher Seifert--van Kampen theorem says
that $U\mapsto\Pi U$ is a cosheaf: an open cover presents its union as the
homotopy colimit of its Cech diagram after applying $\Pi$. Its precise
forms are \citet[Proposition A.3.2, Remark A.3.8, and Lemma A.3.10]{DGLurieHA}.
The last lemma supplies hyperdescent. Pulling an open cover or open
hypercover of $U\subseteq Y$ back to $X$ gives one of $X\times_YU$.
Consequently the same assertions hold for $C_Y(X)$. A map $u$ belongs
to the specified class exactly when $C_Y(u)$ is an objectwise equivalence.
The universal property of localization therefore supplies (D.1).
For the empty open the value is the initial space, as required for a
cosheaf; no nonemptiness or local connectedness assumption was used.
\end{proof}

\begin{claim}{D.2}{Proved}{A complete realization theorem for finite Alexandrov bases}
Let $P$ be a finite poset with the topology of downsets. There are
equivalences
\[
 (\mathrm{Top}/P)[(W^{\mathrm{opn}}_P)^{-1}]
   \simeq\operatorname{Fun}(P,\mathcal S)
   \simeq\operatorname{CoShv}(P;\mathcal S).
\tag{D.2}
\]
The first equivalence sends $f:X\to P$ to the covariant diagram
$p\mapsto\Pi f^{-1}(P_{\leq p})$. It includes arbitrary, possibly empty
fibres and non-Hausdorff total spaces. For a discrete, possibly infinite
base $P$ the same statement holds by decomposition into its open fibres.
\end{claim}
\begin{proof}
We first work before localization. Define
\[
 D:\mathrm{Top}/P\longrightarrow\mathrm{Top}^P,
 \qquad D(X)(p)=X_{\leq p}=f^{-1}(P_{\leq p}).
\]
For a strict covariant diagram $F:P\to\mathrm{Top}$ define $E(F)$ to
have underlying set $\coprod_{r\in P}F(r)$ and projection $(r,z)\mapsto r$.
Give it the topology generated by the subsets
\[
 V(p,O)=\{(r,z):r\leq p,\ F(r\leq p)(z)\in O\},
 \qquad O\subseteq F(p)\text{ open}.
\tag{D.2a}
\]
The projection is continuous since its inverse image of $P_{\leq p}$
is $V(p,F(p))$. A natural transformation $D(X)\to F$ gives the map
$x\mapsto(f(x),\eta_{f(x)}(x))$ to $E(F)$. The inverse image of (D.2a)
is $\eta_p^{-1}(O)$, open in $X_{\leq p}$ and therefore in $X$.
Conversely, a map $X\to E(F)$ over $P$ defines on $X_{\leq p}$ the
continuous function obtained by transporting its value in $F(f(x))$
along $f(x)\leq p$. These constructions are mutually inverse and
natural. Thus $D\dashv E$ as an adjunction of ordinary categories.

The counit at $p$ is
\[
 \epsilon_{F,p}:D E(F)(p)\longrightarrow F(p),
 \qquad(r,z)\longmapsto F(r\leq p)(z).
\]
It has the continuous section $s_p(z)=(p,z)$. On $D E(F)(p)$ the
formula $H(z,t)=z$ for $0\leq t<1$, and
$H(z,1)=s_p\epsilon_{F,p}(z)$, gives a strong deformation retraction.
For a generator $V(q,O)$ with $p\leq q$, membership is unchanged during
this homotopy, by composition in $F$. For a generator with $p\nleq q$,
its inverse image is $(V(q,O)\cap D E(F)(p))\times[0,1)$.
Both sets are open in the product. This verifies continuity even though
the deformation need not be over $P$. Hence $\epsilon_F$ is a pointwise
weak equivalence. The triangle identity
$\epsilon_{D X}\circ D(\eta_X)=1_{D X}$ shows that the unit is a weak
equivalence on every principal open.

For a finite $P$, equivalence on principal opens implies equivalence on
all opens. To see this, use induction on a downset $U$. Choose a maximal
element $p$ of $U$ and write
$U=(U\setminus\{p\})\cup P_{\leq p}$; their intersection is
$P_{<p}$. All three inverse images are open. The open-cover homotopy
pushout theorem of D.1 and induction give the result for $U$. This also
proves that $E$ sends pointwise weak equivalences to
$W^{\mathrm{opn}}_P$: use naturality of its counit, then the induction.
Thus the relative adjunction has unit and counit in the specified
weak-equivalence classes and induces an equivalence of localizations.

The pointwise singular--realization adjunction identifies the localization
of $\mathrm{Top}^P$ with that of $\mathrm{sSet}^P$; its unit and counit
are pointwise weak homotopy equivalences. Strict diagrams of simplicial
sets with their projective weak equivalences present
$\operatorname{Fun}(P,\mathcal S)$ by
\citet[Proposition 4.2.4.4 and Remark 4.2.4.5]{DGLurieHTT}.
This states the exact rectification input; it is not inferred merely
from $\mathrm{Top}[W^{-1}]\simeq\mathcal S$.

Finally, restriction of a cosheaf to principal opens gives a diagram
$P\to\mathcal S$. Its inverse is
$F\mapsto(U\mapsto\operatorname*{colim}_{p\in U}F(p))$.
The same finite-open induction proves codescent and proves that the
unit and counit are equivalences. For discrete $P$, a space over $P$
is the disjoint union of its open fibres, and both sides of (D.2) are
$\prod_{p\in P}\mathcal S$. The empty base gives the terminal category.
\end{proof}

\begin{claim}{D.3}{Refuted}{The two variances are not automatically equivalent unstably}
For finite Alexandrov $P$ the sheaf theory is
$\operatorname{Shv}(P;\mathcal S)\simeq\operatorname{Fun}(P^{\mathrm{op}},\mathcal S)$,
whereas D.2 gives the covariant diagram category for open inverse-image
tests. These categories need not be equivalent. In particular an
unstable sheaf--cosheaf duality for all topological spaces is false.
\end{claim}
\begin{proof}
For a finite poset, a sheaf is recovered from its values on principal
opens by
$F(U)=\lim_{p\in U^{\mathrm{op}}}F(P_{\leq p})$.
Finite gluing, with the induction in D.2 reversed, proves the formula
and its inverse. Take $P=\{b,a,c\}$ with precisely the nonidentity
relations $b<a$ and $b<c$. This poset is not isomorphic to its opposite:
it has one minimal and two maximal elements.

An equivalence between the two space-valued diagram categories would
identify their completely compact objects, meaning objects $K$ for
which $\operatorname{Map}(K,-)$ preserves every small colimit. In a
presheaf category these are exactly retracts of representables. Indeed
write $K$ as its canonical colimit of representables. Preservation of
this colimit makes $1_K$ factor through one representable; this exhibits
$K$ as its retract. Conversely the property is preserved by retracts
and holds for representables. In a poset there are no nontrivial
idempotent endomorphisms, so its category of such retracts is the poset
itself. The two categories would therefore imply $P\simeq P^{\mathrm{op}}$,
which is impossible. The six candidate bijections and their order tests
are exhausted in \texttt{code/dg\_finite\_checks.py}.
This counterexample concerns a finite non-Hausdorff space. It is not
presented as a counterexample under the stronger Hausdorff hypothesis.
\end{proof}

\begin{claim}{D.4}{Refuted}{Proper maps do not preserve open-test equivalences}
Even for locally compact Hausdorff spaces, strict pullback along a
proper map need not preserve $W^{\mathrm{opn}}$. Thus geometric proper
base change for sheaves cannot be read as this assertion about spaces
over a base.
\end{claim}
\begin{proof}
Set $Y=[0,1]$, let $j:(0,1]\hookrightarrow Y$ be a map over $Y$, and
let $i:\{0\}\hookrightarrow Y$. The map $i$ is proper. For any open
$U\subseteq Y$, every connected component of $U$ is an interval. The
component containing $0$, if there is one, stays nonempty and
contractible after $0$ is deleted, and all other components are
unchanged. Consequently $U\cap(0,1]\to U$ induces a bijection on
components and isomorphisms on all based homotopy groups. Thus
$j\in W^{\mathrm{opn}}_Y$. Its strict pullback along $i$ is
$\varnothing\to\{0\}$, which is not a weak homotopy equivalence.
All maps and spaces in this example have been specified; the failure
is on $\pi_0$, not on a choice of derived replacement. Testing instead
on all closed subsets would detect $\{0\}$ and would not invert $j$.
That alternative is therefore a different recipe.
\end{proof}

\begin{claim}{D.5}{Proved}{Local product maps preserve the open-test class}
Suppose $g:Z\to Y$ is locally a projection: each point of $Z$ has an
open neighbourhood isomorphic over an open $V\subseteq Y$ to
$V\times K\to V$, with $K$ a topological space. Then strict pullback
along $g$ sends $W^{\mathrm{opn}}_Y$ to $W^{\mathrm{opn}}_Z$.
\end{claim}
\begin{proof}
Let $u:X\to X'$ be in $W^{\mathrm{opn}}_Y$. On a product rectangle
$V'\times K'\subseteq V\times K$, with both factors open, its
pullback is
\[
 (X\times_YV')\times K'\longrightarrow
 (X'\times_YV')\times K'.
\]
The singular functor preserves products, and products of spaces
preserve weak equivalences, including the case $K'=\varnothing$.
Hence this map is a weak equivalence. Such rectangles form a basis
locally on $Z$. On the basis and its intersection refinements the two
inverse-image homotopy-type functors agree. Hyperdescent in D.1, or
equivalently the basis form of the open-cover homotopy-colimit theorem,
then proves their agreement on every open of $Z$. The argument uses
no compactness or contractibility of $K$.
\end{proof}

\begin{claim}{D.6}{Known}{Proper base change is a coefficient theorem}
Consider a cartesian square of locally compact Hausdorff spaces
\[
\begin{array}{ccc}
 X'&\xrightarrow{g'}&X\\[-2pt]
 {\scriptstyle f'}\downarrow&&\downarrow{\scriptstyle f}\\[-2pt]
 Y'&\xrightarrow{g}&Y.
\end{array}
\]
If $f$ is proper, the canonical mate
$g^*f_*\to f'_*g'^*$ is an equivalence on sheaves of spaces.
It is also an equivalence on sheaves with values in a stable
bicomplete infinity-category. These are ordinary infinity-sheaves;
hypercompletion is not silently substituted.
\end{claim}
\begin{proof}
The space-valued statement is exactly
\citet[Corollary 7.3.1.18]{DGLurieHTT}. Its proof first establishes
properness of the induced geometric morphism and then the appropriate
pullback of infinity-topoi; the displayed comparison is the
adjointability mate in that theorem. For stable coefficients the
extension and its hypotheses are given by
\citet[Section 6.1 and Proposition 6.9]{DGVolpe2025}, using $f_!=f_*$
for proper $f$. This identifies the specified map, not merely two
abstractly equivalent functors. D.4 explains why the theorem cannot
be transferred to strict fibres of arbitrary spaces over $Y$.
\end{proof}

\begin{claim}{D.7}{Known}{A precise smooth topological comparison}
A space $K$ is essential when the constant-sheaf functor
$\mathcal S\to\operatorname{Shv}(K;\mathcal S)$ has a left adjoint.
A shape submersion $g:Z\to Y$ is locally, on its source, a projection
$V\times K\to V$ with $V$ open in $Y$ and $K$ essential.
Assume either $Y$ is locally compact or all these $K$ are locally
compact. Then for any presentable coefficient infinity-category
$\mathcal C$ there is an adjunction
$g_\sharp\dashv g^*$ on sheaves with these coefficients. In a
cartesian square as in D.6 with $g$ a shape submersion and with the
same local compactness condition after pullback, the canonical maps
\[
 g'_\sharp f'^*\longrightarrow f^*g_\sharp,
 \qquad g^*f_*\longrightarrow f'_*g'^*
\]
are equivalences.
\end{claim}
\begin{proof}
These are Definition 3.23, the hypothesis immediately preceding
Lemma 3.26, and Lemmas 3.26--3.27 of \citet{DGVolpe2025}, with the
square relabelled. Their proof reduces to the indicated product
charts, where essentiality supplies the left adjoint, and then uses
descent. The second map is the comparison between the full right
adjoints of the first. This is a specified geometric sufficient
condition, not a definition of ``locally aspherical'' by the
homotopy types of individual fibres. D.5 additionally establishes
preservation of open-test equivalences for this class.
\end{proof}

\begin{claim}{D.8}{Known}{The stable sheaf--cosheaf equivalence and the shriek operation}
For a locally compact Hausdorff space $Y$ and a stable bicomplete
infinity-category $\mathcal C$, there is an equivalence
\[
 \mathbb V_Y:\operatorname{Shv}(Y;\mathcal C)
 \xrightarrow{\ \simeq\ }\operatorname{CoShv}(Y;\mathcal C),
 \qquad
 \mathbb V_Y(F)(U)=
 \operatorname*{colim}_{K\subseteq U\ {
m compact}}
 \operatorname{fib}\bigl(F(Y)\to F(Y\setminus K)\bigr).
\tag{D.8}
\]
For a continuous map $f:X\to Y$ between such spaces, cosheaf
pushforward $f_+$ is $(f_+G)(U)=G(f^{-1}U)$. The topological operation
with proper support is
$f_!=\mathbb V_Y^{-1}f_+\mathbb V_X$ and has a right adjoint $f^!$.
For proper $f$ it agrees with $f_*$. This does not identify $f_!$ with
the left adjoint $f_\sharp$ of D.7.
\end{claim}
\begin{proof}
Equation (D.8) is \citet[Theorem 5.5.5.1 and Proposition 5.5.5.10]{DGLurieHA},
with the opposite categories in that theorem taken on both sides.
Existence of the indicated adjoints and the proper identification are
established in \citet[Section 6.1]{DGVolpe2025}. They depend on
compact support, absent from the definition of an ordinary Kan
extension on a category. For example, for $a:\mathbb R\to *$ and
constant coefficients $k$, $a_\sharp k\simeq k$ while
$a_!k\simeq k[-1]$ in cohomological grading. The first equality is
homotopy invariance on the contractible line; the second follows by
one-point compactification, whose reduced cohomology is $k$ in
degree one. Thus even on a manifold the two left adjoints have
different right adjoints. Composing (D.8) with coefficient dualization
gives the usual contravariant Verdier duality only on subcategories
where that coefficient dualization is an equivalence.
\end{proof}

\begin{claim}{D.9}{Refuted}{Realization does not turn every Kan extension into a geometric direct image}
Under the least-vertex stratification, categorical properness and
properness of a realized finite-poset map are different. Ambient
topological direct image need not preserve the least-vertex
constructible sheaves and need not equal categorical right Kan
extension.
\end{claim}
\begin{proof}
Let $P=[1]$, so $|NP|=[0,1]$ has strata $[0,1)$ and $\{1\}$.
Let $i:*\to P$ select $0$. Its realization is a proper closed
inclusion, while $i$ is not categorically proper: the arrow $0\to1$
has no outgoing lift. A presheaf on $P$ is an arrow $F_1\to F_0$.
The right Kan extension of a two-point set $M$ from $0$ is the
constant arrow $M\xrightarrow{1}M$, since both comma categories
$*/0$ and $*/1$ are terminal. On the interval, direct image from
the endpoint is the skyscraper: its stalk is $M$ at $0$ and the
terminal space at every other point. It is not locally constant on
$[0,1)$. The categorical sheaf would instead be constant on that
stratum. The finite category, its comma categories, and the value
cardinalities are verified in \texttt{code/dg\_finite\_checks.py}.
\end{proof}

\begin{claim}{D.10}{Proved}{A geometric class on which the categorical comparison is literal}
For finite posets equipped with their least-vertex models, closed
cosieve inclusions and product projections have topological direct
images equal to their categorical right Kan extensions for space,
spectrum, and derived-ring coefficients. A sieve is represented by
the open complement of the complementary cosieve, rather than in
general by the inclusion of its own nerve. These assertions use the
finite-poset exodromy comparison, as verified in G.10 below.
\end{claim}
\begin{proof}
If $i:B\hookrightarrow P$ is a cosieve, $|NB|$ is a closed
subcomplex and is a union of least-vertex strata: a chain whose
minimum belongs to $B$ lies entirely in $B$. Topological closed
direct image is terminal off this subcomplex and agrees with the
given sheaf on it. It preserves constructibility and is right
adjoint to the specified restriction, so uniqueness of right
adjoints identifies it with $i_*$. This also identifies units,
counits, and base-change mates.

If $A=P\setminus B$ is the complementary sieve, let
$O=|NP|\setminus|NB|$. A point of $O$ has positive total barycentric
coordinate on vertices of $A$. Delete the other coordinates and
renormalize; straight-line interpolation with this operation gives
a deformation retraction $O\to|NA|$. Its minimum vertex is unchanged
throughout. Constructible exodromy therefore identifies coefficients
on $O$ with coefficients on $A$, including restriction and its
adjoints. The inclusion $|NA|\hookrightarrow|NP|$ itself need not
be open, as $\{0\}\hookrightarrow[0,1]$ demonstrates.

Finally the canonical homeomorphism
$|N(P\times Q)|\cong|NP|\times|NQ|$ respects the minimum strata.
Projection has compact fibre $|NP|$. Over a stratum of $Q$, the
given sheaf is a family of fixed constructible sheaves on $|NP|$
with locally constant transport; taking its sections leaves that
transport locally constant. This follows as well by the finite
exodromy diagram calculation, since the sections are a finite
homotopy limit. The direct image thus preserves constructibility
and is the unique right adjoint to restriction. D.6 applies to its
geometric pullback squares. For the three stated coefficient
categories the finite exodromy comparison is available, and finite
limits commute with their locally constant transport.
\end{proof}

\begin{claim}{D.11}{Proved}{The realization comparison for every proper map of finite posets}
Let $f:P\to Q$ be a $W_\infty$-proper map of finite posets.
Equivalently, for every $p\in P$ and $q\geq f(p)$ the outgoing
lifting poset
\[
 L(p,q)=\{r\in P:r\geq p,\ f(r)=q\}
\tag{D.11a}
\]
is weakly contractible. With the minimum stratifications, the
topological direct image $|Nf|_*$ preserves constructibility and
identifies with the categorical $f_*$, including its adjunction
data. This holds for spaces, spectra, and the derived category of
modules over any ring. Thus the requested geometric class contains
the realizations of all such proper functors, with the specified
constructibility property, although mere topological properness
does not suffice by D.9.
\end{claim}
\begin{proof}
Write $g=|Nf|$. We must compare the fibres at interior simplex
points as well as at vertices. Fix a point $x\in|NQ|$ whose
positive barycentric support is the chain
$\sigma=(q_0<\cdots<q_n)$, with weights $\lambda_i>0$.
The fibre $Z_x=g^{-1}(x)$ is a finite polytopal complex. Its cells
are indexed by nonempty chains $\kappa$ of $P$ with
$f(\kappa)=\sigma$. The cell indexed by $\kappa$ is the product
of open simplices
\[
 \prod_{i=0}^n\operatorname{int}\Delta^{\kappa\cap P_{q_i}}:
 \quad (t_p)_{p\in\kappa\cap P_{q_i}}
            \text{ has sum }\lambda_i.
\tag{D.11b}
\]
The closure is the corresponding product of closed simplices.
Its faces are precisely the subchains still meeting every level.
Thus the face poset is
\[
 C_\sigma=\{\kappa\subseteq P:\kappa\text{ a chain},\
                          f\kappa=\sigma\},
 \quad\text{ordered by inclusion}.
\]
The restriction of a minimum-constructible sheaf associated to
$F:P^{\mathrm{op}}\to\mathcal C$ has value $F(\min\kappa)$ on
this cell, with its prescribed structure maps. The exit category
of this finite polytopal cell stratification is its face poset:
barycentrically triangulate the cells, use HA A.6.10 on the
triangulation, and coarsen by the maximum face using HPT 3.3.5 and
the last-vertex localization, exactly as in G.10. Hence
\[
 \Gamma(Z_x;F|_{Z_x})
 \simeq\lim_{\kappa\in C_\sigma}F(\min\kappa).
\tag{D.11c}
\]
This uses the coefficient classes in G.10, and applies to the
three classes in the statement.

We now prove the required initiality, rather than assuming an
exit-category equivalence for the whole fibre. The minimum functor
\[
 m:C_\sigma\longrightarrow(P_{q_0})^{\mathrm{op}}
\]
is homotopy initial. For $p\in P_{q_0}$, its comma category $m/p$
is the poset of chains in $C_\sigma$ whose every vertex is at least
$p$. We prove that this poset is weakly contractible by induction
on $n$. Add $p$ to every chain, then delete every other vertex at
level $q_0$. The two operations are connected to the identity by
the natural inclusions
\[
 \kappa\ \subseteq\ \kappa\cup\{p\}\
 \supseteq\ \{p\}\cup\{r\in\kappa:f(r)>q_0\}.
\]
They are retractions, up to these natural transformations, onto
the subposet with exactly $p$ at its first level. For $n=0$ this
subposet is the one-object poset $\{\{p\}\}$.

For $n>0$, delete $p$ from this subposet and send each remaining
chain to its minimum at level $q_1$. This gives a functor to
$L(p,q_1)^{\mathrm{op}}$. Its comma category over an object
$r\in L(p,q_1)$ consists of the chains over
$(q_1<\cdots<q_n)$ all of whose vertices are at least $r$.
It is weakly contractible by the induction hypothesis: all the
further lifting posets $L(r,q_i)$ and those above their elements
are exactly the posets covered by (D.11a). Quillen's Theorem A
therefore makes this minimum functor a weak equivalence of nerves.
Its target is weakly contractible by (D.11a), proving the induction.
The same Theorem A, now applied to $m$, proves homotopy initiality.
Consequently (D.11c) is canonically
\[
 \Gamma(Z_x;F|_{Z_x})\simeq
             \lim_{(P_{q_0})^{\mathrm{op}}}F.
\tag{D.11d}
\]
If the fibre $P_{q_0}$ is empty both indexing categories are empty
and both limits are terminal; the argument covers that case too.

Categorical right Kan extension has the pointwise value
\[
 (f_*F)(q_0)=\lim_{(P_{\leq q_0})^{\mathrm{op}}}F,
 \qquad P_{\leq q_0}=\{p:f(p)\leq q_0\}.
\]
The inclusion $(P_{q_0})^{\mathrm{op}}\to
(P_{\leq q_0})^{\mathrm{op}}$ is homotopy initial: its comma
category at $p$ is $L(p,q_0)^{\mathrm{op}}$, contractible by
(D.11a). Thus this Kan value also equals the right side of (D.11d).

It remains to identify the comparison and detect it globally.
Let $E_P,E_Q$ denote the fully faithful exodromy functors of G.10.
Minimum labels commute with $f$, so
$g^*E_Q\simeq E_Pf^*$. The counit gives the canonical mate
\[
 E_Qf_*F\longrightarrow g_*E_PF.
\tag{D.11e}
\]
The map $g$ is proper because it is a map of compact Hausdorff
finite polyhedra. Proper base change identifies the target stalk
at $x$ with (D.11c). The source stalk is $(f_*F)(q_0)$.
The stalk of (D.11e) is precisely restriction of the limit cone
along
$C_\sigma\xrightarrow{m}(P_{q_0})^{\mathrm{op}}
\hookrightarrow(P_{\leq q_0})^{\mathrm{op}}$:
the counit at a cell with minimum $p$ is evaluation at that same
object $p$ of the Kan indexing category. Both functors are
homotopy initial, so this specified map is an equivalence.

Finally sheaves of spaces on a finite polyhedron are hypercomplete:
finite covering dimension and paracompactness give finite homotopy
dimension by \citet[Corollary 7.2.3.7]{DGLurieHTT}, and
\citet[Corollary 7.2.1.12]{DGLurieHTT}
then gives hypercompleteness. Stalks therefore detect equivalences.
For spectra and modules the same conclusion follows by applying
maps from shifts of the sphere or the free rank-one module;
these conservative functors commute with filtered colimits and
limits, so they commute with stalks and sheaf evaluation.
Thus (D.11e) is an equivalence of sheaves. Its source is
constructible, proving preservation, and its construction as a
mate gives the asserted compatibility with units and counits.
\end{proof}

\begin{claim}{D.12}{Open}{The unrestricted topological realization and its proposed dual}
For a locally compact Hausdorff base $Y$, is the functor $C_Y$ in
(D.1), with all topological total spaces allowed, an equivalence?
If so, what precise localization of spaces over the same $Y$
presents $\operatorname{Shv}(Y;\mathcal S)$ and is compatible with
the proposed dual class of base-change maps? Neither statement is
established in this audit for a general $Y$; the first unproved
test base in our argument is $\mathbb R$.
\end{claim}
The construction in D.2 needs a smallest open neighbourhood for each
point; a real line has none. D.1 gives a functor and detects exactly
the original class on strict maps, but does not prove essential
surjectivity or full faithfulness after localization. Testing every
closed subset defines a legitimate different weak-equivalence
class, as D.4 shows, but this fact alone supplies neither descent
nor an equivalence with sheaves. We found no primary theorem in the
searched sources establishing that unrestricted representation
statement. Lurie's stable sheaf--cosheaf theorem D.8 makes no such
assertion about $\mathrm{Top}/Y$ and cannot fill this gap.

\begin{claim}{D.13}{Proved}{A specified stable cosheaf counterpart for smooth finite-poset maps}
There is a precise geometric dual on finite coefficients which uses
the maximum stratification and Verdier transport. Let $\mathcal C$
be spectra or modules over a commutative ring, let $D$ be the
coefficient duality on finite spectra or perfect modules, and let
$E_P^{\max}:\operatorname{Fun}(P,\mathcal C)\to
\operatorname{Shv}(|NP|;\mathcal C)$ be G.10 for $P^{\mathrm{op}}$.
Define the contravariant fully faithful functor
\[
 \Phi_P(F)=\mathbb V_{|NP|}\,E_P^{\max}(DF),
 \qquad F\in\operatorname{Fun}(P^{\mathrm{op}},\mathcal C^{\mathrm{perf}}).
\tag{D.13a}
\]
If $f:P\to Q$ is categorically smooth, then its realization $g$
satisfies
\[
 \Phi_Q(f_!F)\simeq g_+\Phi_P(F).
\tag{D.13b}
\]
Thus ordinary cosheaf pushforward realizes this specified dual of
left Kan extension. The cosheaf subcategories in (D.13a) are
defined as Verdier images of maximum-constructible perfect sheaves;
they are not being identified with literal minimum-constructible
cosheaves. For a noncommutative ring the target coefficients are
modules over its opposite ring.
\end{claim}
\begin{proof}
Smoothness of $f$ is properness of $f^{\mathrm{op}}$.
Apply D.11 to $f^{\mathrm{op}}:P^{\mathrm{op}}\to Q^{\mathrm{op}}$.
Its minimum stratification is the maximum stratification on the
same underlying polyhedron. Coefficient duality on finite diagrams
gives
$D(f_!F)\simeq(f^{\mathrm{op}})_*(DF)$: all comma categories are
finite direct, so the relevant limits and colimits remain perfect,
and the duality in G.6 exchanges them. Consequently D.11 gives
\[
 E_Q^{\max}D(f_!F)\simeq g_*E_P^{\max}(DF).
\]
The map $g$ is proper, and D.8 identifies
$\mathbb V_{|NQ|}g_*\simeq g_+\mathbb V_{|NP|}$.
Applying this equivalence proves (D.13b), including the canonical
maps obtained from these adjunctions. Full faithfulness follows
from coefficient biduality, G.10, and the equivalence D.8.
The description of the image is part of the construction, so this
does not make an unproved constructibility assertion about
cosheaves for a coarser minimum stratification; G.11 illustrates why
that assertion would require a separate argument.
\end{proof}
